% Standalone R2 dossier for the elementary linear-test and Gaussian-prior GLM baselines.
\documentclass[../main.tex]{subfiles}
\begin{document}

% === SOLUTION HEADER (proof provenance) =====================================
% ledger-nodes: lem:linear-test-lower; thm:glm-fi
% refines: rem:rayleigh-lower-only; thm:glm-fi
% bounded_by: none
% evidence_run: none
% checked_by: agent
% author: /root/ab_legacy_author
% reviewer: /root/kls_core_author
% review: research/reviews/2026-08-25-ab-inline-r2-audit.md
% date: 2026-08-25
% ============================================================================

\section*{Solution: linear-test and Gaussian-prior GLM baselines}
\label{sec:sol-glm-linear-baselines}

\paragraph{Normalization and imported input.}
We use Definitions~\ref{def:poincare}--\ref{def:tci}: in particular,
\(\Ent_\pi(f^2)\le 2C\int\norm{\nabla f}^2\dd\pi\) and
\(W_2^2(\eta,\pi)\le2C\KL(\eta\|\pi)\).  The only imported results in the
GLM proof are the following standard theorems, in exactly this normalization.

\begin{enumerate}[leftmargin=1.7em]
  \item If a smooth probability density is proportional to \(e^{-U}\) and
  \(\Hess U\succeq mI\), \(m>0\), the Bakry--\'Emery criterion gives
  \(\CP\le\CLS\le m^{-1}\); this is Theorem~\ref{thm:bakry-emery}
  \cite{BakryEmery1985,BakryGentilLedoux2014}.
  \item The Otto--Villani implication gives \(\CTCI\le\CLS\) under the same
  normalization; this is Theorem~\ref{thm:otto-villani}
  \cite{OttoVillani2000,Villani2008}.
\end{enumerate}

The approximation argument below is included so that convexity in
Family~\ref{fam:glm-posterior} is not silently strengthened to \(C^2\)-smoothness.

\subsection*{1. The universal linear-test lower bound}

\begin{lemma}[Ledger node \texttt{lem:linear-test-lower}]
\label{lem:sol-linear-test-lower}
Let \(\pi\) be a probability measure on \(\R^D\) with finite covariance matrix and let
\(\CP(\pi)\in[0,+\infty]\) be its optimal Poincar\'e constant.  Every admissible nonconstant
test function \(f\) satisfies
\[
 \frac{\Var_\pi(f)}{\int\norm{\nabla f}^2\dd\pi}\le \CP(\pi),
\]
with the usual extended-value convention.  In particular,
\[
 \CP(\pi)\ \ge\ \lmax(\Cov_\pi).
\]
This is the claim recorded at Remark~\ref{rem:rayleigh-lower-only}.
\end{lemma}

\begin{proof}
The first assertion is the defining Poincar\'e inequality rearranged; equivalently, the optimal
constant is the supremum of its Rayleigh quotients.  For \(v\in\R^D\), take
\(f_v(x)=\inner{v}{x}\).  Then
\[
 \Var_\pi(f_v)=v^T\Cov_\pi v,
 \qquad
 \int\norm{\nabla f_v}^2\dd\pi=\norm v^2.
\]
Consequently \(v^T\Cov_\pi v\le \CP(\pi)\norm v^2\) for every \(v\).  Taking the supremum
over unit vectors gives \(\lmax(\Cov_\pi)\le\CP(\pi)\).  If \(\CP(\pi)=+\infty\), the same
conclusion is immediate.
\end{proof}

\subsection*{2. The data-free Gaussian-prior GLM bound}

\begin{theorem}[Ledger node \texttt{thm:glm-fi}]
\label{thm:sol-glm-fi}
Let \(\Sigmaz\) be a positive-definite \(D\times D\) matrix, let
\(x_1,\ldots,x_n\in\R^D\), and let each scalar loss
\(t\mapsto\ell_i(t)=\ell(y_i,t)\) be finite and convex on \(\R\).  Define the
Gaussian-prior GLM posterior of Family~\ref{fam:glm-posterior} by
\[
 \pi(\dd\theta)=Z^{-1}\exp\!\left[-\frac12\theta^T\Sigmaz^{-1}\theta
              -\sum_{i=1}^n\ell_i(x_i^T\theta)\right]\dd\theta .
\]
Then \(0<Z<\infty\), and
\[
 \CP(\pi)\le\CLS(\pi)\le\lmax(\Sigmaz),
 \qquad
 \CTCI(\pi)\le\lmax(\Sigmaz).
\]
The displayed certificate depends on the data only through the fact that the likelihood loss is
convex and has no explicit dimension factor; it is therefore the data-free, dimension-free
baseline asserted in Theorem~\ref{thm:glm-fi}.
\end{theorem}

\begin{proof}
Put \(m=\lmin(\Sigmaz^{-1})=\lmax(\Sigmaz)^{-1}>0\).  First suppose that every
\(\ell_i\) is \(C^2\).  The posterior potential
\[
 U(\theta)=\frac12\theta^T\Sigmaz^{-1}\theta+
           \sum_{i=1}^n\ell_i(x_i^T\theta)
\]
satisfies
\[
 \Hess U(\theta)=\Sigmaz^{-1}+
 \sum_{i=1}^n\ell_i''(x_i^T\theta)x_ix_i^T
 \succeq\Sigmaz^{-1}\succeq mI_D,
\]
because convexity gives \(\ell_i''\ge0\).  Bakry--\'Emery therefore yields
\(\CP(\pi)\le\CLS(\pi)\le m^{-1}=\lmax(\Sigmaz)\), and Otto--Villani yields
\(\CTCI(\pi)\le\CLS(\pi)\).

We now remove the smoothness assumption.  A finite convex function on \(\R\) is continuous and
has a supporting affine function: choose \(s_i\in\partial\ell_i(0)\), so
\[
 \ell_i(t)\ge \ell_i(0)+s_it\qquad(t\in\R).
\]
It follows that \(U\) is bounded below by a positive-definite quadratic plus an affine function.
Thus \(e^{-U}\) has a Gaussian integrable majorant, which proves \(0<Z<\infty\).

Let \(\rho_\delta\) be a smooth, compactly supported, even mollifier on \(\R\), and set
\(\ell_{i,\delta}=\rho_\delta*\ell_i\).  These functions are smooth and convex, converge
locally uniformly to \(\ell_i\), and, because the mollifier has mean zero, obey the same affine
lower bound
\[
 \ell_{i,\delta}(t)\ge \ell_i(0)+s_it.
\]
Let \(\pi_\delta\) be the corresponding posterior.  The smooth argument gives, uniformly in
\(\delta\),
\[
 \Var_{\pi_\delta}(f)\le m^{-1}\int\norm{\nabla f}^2\dd\pi_\delta,
 \qquad
 \Ent_{\pi_\delta}(f^2)\le2m^{-1}\int\norm{\nabla f}^2\dd\pi_\delta.
\]
The common affine lower bounds provide one Gaussian integrable majorant for the unnormalised
densities.  Dominated convergence therefore gives convergence of the normalizing constants and
of all terms in these two displays for bounded smooth compactly supported \(f\).  Passing to the
limit proves the same Poincar\'e and log-Sobolev inequalities for \(\pi\) on that core; truncation
and closure extend them to their natural Sobolev domains.  Applying the imported Otto--Villani
theorem directly to this limiting log-Sobolev inequality proves the stated \(T_2\) bound.  This
also shows explicitly that no differentiability of the original convex losses is being assumed.
\end{proof}

\paragraph{Scope.}
The theorem asserts only the prior-curvature certificate.  It does not claim that the optimal
constants are data-independent or equal to \(\lmax(\Sigmaz)\), and it does not supply the
data-informed sharpening sought by A1.

\end{document}
